% !TeX program = pdflatex
% =============================================================================
% Arbeitsblatt 1: Elektrische Feldlinien zeichnen
% UZH Praktikum I -- Sara Romer Urban, Kantonsschule im Lee, HS2026
% Klasse: 1f (02.09.2026), aus PLAN_Dokumentengenerierung.md.
%
% Bewusst schlank gehalten (Wunsch: "ohne viel Text, es soll nur ums
% E-Feld zeichnen gehen"): keine Lernziele-/Einleitungsbox (wie bereits bei
% kraefte-addieren-zerlegen/magnetfelder in diesem Ordner -- Praktikum I
% startet direkt mit Inhalt). Coulombkraft/Anziehung-Abstossung sowie eine
% erste Einfuehrung ins elektrische Feld sind laut Vorwissen/Elektrostatik
% (elektrodynamik/Vorwissen/Elektrostatik/Arbeitsblatt1+2) bereits bekannt
% -- deshalb hier nur eine kurze Repetition, keine erneute Herleitung.
%
% Die vier Zeichenuebungen (Einzelladung +/-, Dipol, Plattenkondensator)
% zeigen als Aufgabe nur die tikz-Skizze der Ladungen/Platten (leeres
% Feld zum Einzeichnen); die fertigen Feldlinienbilder liegen als Bilder/
% bereits vor und erscheinen erst in der Loesung (Konvention "SuS selbst
% ergaenzen lassen", institutionen/UZH-Praktikum-I/CLAUDE.md).
%
% Numerierung "Arbeitsblatt1": elektrizitaet/Vorwissen (der fuer 1f
% korrekte Bereichsordner) enthaelt noch keine nummerierten Dokumente
% Sara Romers -- nur die scheinbar dazugehoerigen Arbeitsblatt1-3 liegen
% (vermutlich falsch einsortiert) unter elektrodynamik/Vorwissen/
% Elektrostatik/. Falls das die tatsaechliche 1f-Nummerierung fortsetzt,
% mit Sara Romer abgleichen und Dateiname ggf. anpassen.
% =============================================================================
\documentclass[addpoints,11pt,a4paper]{exam}

\usepackage[T1]{fontenc}
\usepackage[utf8]{inputenc}
\usepackage[ngerman]{babel}
\usepackage{lmodern}
\usepackage[a4paper,margin=2.2cm,headheight=14pt]{geometry}
\usepackage{amsmath}
\usepackage{siunitx}
% Hausstil: zusammengesetzte Einheiten immer als echter Bruch (\frac), nicht
% als Inline-Schraegstrich oder \tfrac -- gilt global fuer jedes \si/\SI.
\sisetup{per-mode=fraction,fraction-command=\frac}
\usepackage{array,booktabs,tabularx,multirow}
\usepackage{enumitem}
\usepackage{xcolor}
\usepackage{colortbl}
\usepackage{tikz}
\usepackage{physics}
\usepackage[most]{tcolorbox}
\usepackage{lastpage}
\usepackage{graphicx}
\usepackage{hyperref}

\usetikzlibrary{arrows.meta,calc,angles,quotes}

\usepackage{setspace}
\usepackage{needspace}
\setlength{\parindent}{0pt}
\setlength{\parskip}{0.6em}
\setstretch{1.08}
\setlist[itemize]{itemsep=0.4em}

% -----------------------------
% Farben (Corporate-Farbschema Physik KFR)
% -----------------------------
\definecolor{titleblue}{HTML}{183A6B}
\definecolor{accentblue}{HTML}{2E6F9E}
\definecolor{lightblue}{HTML}{EAF4FB}
\definecolor{lightgray}{HTML}{F4F6F8}
\definecolor{darkgray}{HTML}{4A4A4A}
\definecolor{accentorange}{HTML}{D9720C}
\definecolor{accentpurple}{HTML}{7B2D8E}
\definecolor{lightorange}{HTML}{FCEEE0}
\definecolor{lightpurple}{HTML}{F3EAF7}

% Links (hyperref): farblich ans Corporate-Farbschema angeglichen.
\hypersetup{
  colorlinks=true,
  urlcolor=accentblue,
  linkcolor=accentblue,
  pdftitle={Arbeitsblatt 1: Elektrische Feldlinien zeichnen},
  pdfauthor={Loris Keller}
}

% -----------------------------
% Spaltentypen
% -----------------------------
\newcolumntype{Y}{>{\centering\arraybackslash}X}
\newcolumntype{P}[1]{>{\raggedright\arraybackslash}p{#1}}

% -----------------------------
% Boxen
% -----------------------------
\tcbset{before skip=10pt, after skip=10pt}
\newtcolorbox{titlebox}{
  enhanced,
  colback=titleblue,
  colframe=titleblue,
  boxrule=0pt,
  arc=3mm,
  left=3mm,right=3mm,top=3mm,bottom=3mm
}

\newlength{\aufgabenindent}
\settowidth{\aufgabenindent}{10.\hskip\labelsep}

% Praktikum I: Lernziele/Einleitung/Theorie/Aufgaben werden NICHT farbig
% gerahmt (nur Formel-, Hinweis- und Antwortbox bleiben Boxen) -- siehe
% institutionen/UZH-Praktikum-I/CLAUDE.md, Abschnitt "Boxen-Layout".
\newenvironment{infobox}{%
  \par\addvspace{10pt}\noindent
}{%
  \par\addvspace{10pt}
}

\newenvironment{theoriebox}[1][Theorie]{%
  \par\addvspace{10pt}\noindent\textbf{#1}\par\smallskip\noindent
}{%
  \par\addvspace{10pt}
}

\newtcolorbox{taskbox}[1][]{
  enhanced, breakable,
  colback=white, colframe=white, boxrule=0pt,
  left=0mm,right=0mm,top=0mm,bottom=0mm,
  before skip=10pt, after skip=10pt,
  before upper={%
    \def\tmpboxtitle{#1}%
    \ifx\tmpboxtitle\empty\else\noindent\textbf{#1}\par\smallskip\fi
  }
}

\newtcolorbox{groupbox}[1][Gruppenarbeit]{
  enhanced, breakable,
  colback=white, colframe=white, boxrule=0pt,
  left=0mm,right=0mm,top=0mm,bottom=0mm,
  before skip=10pt, after skip=10pt,
  before upper={\noindent\textbf{#1}\par\smallskip}
}

\newenvironment{beispielbox}[1][Beispiel]{%
  \par\addvspace{10pt}\noindent\textbf{#1}\par\smallskip\noindent
}{%
  \par\addvspace{10pt}
}

\newtcolorbox{hinweisbox}[1][Hinweis]{
  enhanced,
  breakable,
  colback=lightorange,
  colframe=accentorange,
  title={#1},
  fonttitle=\bfseries,
  boxrule=0.7pt,
  arc=2mm,
  left=5mm,right=5mm,top=2mm,bottom=2mm
}

\newtcolorbox{formelbox}[1][Formel]{
  enhanced,
  breakable,
  colback=lightpurple,
  colframe=accentpurple,
  title={#1},
  fonttitle=\bfseries,
  boxrule=0.9pt,
  arc=2mm,
  left=5mm,right=5mm,top=2mm,bottom=2mm
}

\newtcolorbox{answerbox}{
  breakable,
  colback=white,
  colframe=gray!55,
  boxrule=0.5pt,
  arc=1.5mm,
  left=2mm,right=2mm,top=2mm,bottom=2mm
}

% Bild-Helfer: zentriertes Bild mit spuerbarem Abstand zum umgebenden Text.
% #1 = Breite, #2 = Dateipfad.
\newcommand{\Bild}[2]{%
  \par\vspace{0.6cm}
  \begin{center}
    \includegraphics[width=#1]{#2}
  \end{center}
  \vspace{0.6cm}
}

% Ladungssymbole fuer die Zeichenuebungen (rot/blau, wie in den
% Loesungsbildern in Bilder/).
\tikzset{
  posladung/.style={circle,fill=red,draw=black,thick,minimum size=8mm,
    font=\bfseries\Large\color{white}},
  negladung/.style={circle,fill=blue,draw=black,thick,minimum size=8mm,
    font=\bfseries\Large\color{white}}
}

% --- Loesungen ein-/ausblenden -----------------------------------------------
\noprintanswers
%\printanswers

\pointformat{}
\renewcommand{\solutiontitle}{\noindent\color{titleblue}\textbf{L\"osung:}\enspace}

\pagestyle{headandfoot}
\cfoot{\textcolor{darkgray}{Seite \thepage\ von \pageref{LastPage}}}

\newcommand{\Kopfzeile}[4]{%
  \begin{titlebox}
    {\color{white}\sffamily\bfseries\Large #4}
  \end{titlebox}
  \vspace{0.5em}
  \noindent\textbf{Klasse:}~#1 \hfill \textbf{Datum:}~#2 \hfill \textbf{Thema:}~#3
    \hfill \textbf{Lehrperson:}~Loris Keller
  \vspace{0.8em}
  \lhead{\textcolor{darkgray}{#3}}
  \rhead{\textcolor{darkgray}{#4}}
}

\newlength{\antwortraumlen}
\newcommand{\Antwortraum}[1]{%
  \pgfmathsetlength{\antwortraumlen}{#1*2.5cm}%
  \begin{answerbox}
    \rule{0pt}{\antwortraumlen}%
  \end{answerbox}%
}

% Werte fuer Kopfzeile zentral setzen
\newcommand{\Klasse}{1f}
\newcommand{\Datum}{02.09.2026}
\newcommand{\Thema}{Elektrisches Feld}
\newcommand{\ArbeitsblattTitel}{Arbeitsblatt 1: Elektrische Feldlinien zeichnen}

\begin{document}

\Kopfzeile{\Klasse}{\Datum}{\Thema}{\ArbeitsblattTitel}

% =============================================================================
% Repetition: Anziehung/Abstossung, Influenz
% =============================================================================
\textbf{Repetition}

Sie wissen bereits: Gleichnamige Ladungen stossen sich ab, ungleichnamige
ziehen sich an. Diese Kraft wirkt als Fernwirkung, ganz ohne Berührung:
Über das elektrische Feld spürt ein Körper die Anwesenheit einer Ladung
auch aus einiger Entfernung. Nähert sich eine Ladung sogar einem
neutralen Leiter, verschieben sich dort dessen Ladungsträger. Das nennt
man \textbf{\emph{Influenz}}. Genau diese Wirkung stellen wir jetzt
bildlich dar: mit \textbf{\emph{elektrischen Feldlinien}}.

% =============================================================================
% Theorie: Regeln zum Zeichnen elektrischer Feldlinien
% =============================================================================
\begin{theoriebox}[Elektrische Feldlinien zeichnen]
\textbf{\emph{Regeln zum Zeichnen elektrischer Feldlinien:}}
\begin{itemize}[leftmargin=1.2cm,itemsep=0.2em]
  \item Feldlinien beginnen bei \textbf{\emph{positiven}} Ladungen und enden bei \textbf{\emph{negativen}} Ladungen. Bei nur einer einzelnen Ladung enden sie im Unendlichen.
  \item Die Richtung einer Feldlinie an einem Punkt ist die Richtung, in die eine dort platzierte kleine \textbf{\emph{positive Probeladung}} gedrückt würde.
  \item Feldlinien kreuzen oder berühren sich nie.
  \item Je dichter die Feldlinien gezeichnet sind, desto \textbf{\emph{stärker}} ist das Feld an dieser Stelle.
  \item Auf der Oberfläche eines geladenen Leiters (z.\,B. einer Platte) stehen Feldlinien immer \textbf{\emph{senkrecht}}.
\end{itemize}

\textbf{\emph{Homogenes Feld:}} Verlaufen die Feldlinien überall parallel
und gleich dicht, ist das Feld an jeder Stelle gleich stark und gleich
gerichtet. Man nennt es dann \textbf{\emph{homogen}}. Ändern sich Richtung
oder Dichte von Ort zu Ort, ist das Feld \textbf{\emph{inhomogen}}.

\begin{hinweisbox}[Ihre Eselsbrücke bei jeder Aufgabe]
Fragen Sie sich an jedem Punkt: \textbf{\emph{Was würde mit einer kleinen
positiven Probeladung genau hier passieren?}} In diese Richtung zeigt die
Feldlinie an dieser Stelle.
\end{hinweisbox}
\end{theoriebox}

% =============================================================================
% Aufgabe 1: Einzelne positive Ladung
% =============================================================================
\needspace{9.5cm}
\begin{questions}
\begin{taskbox}[Aufgabe 1: Einzelne positive Ladung]
\question[4]
\begin{parts}
  \part[3] Zeichnen Sie die Feldlinien um die abgebildete positive
  Ladung ein.
  \begin{answerbox}
  \vspace{2.1cm}
  \begin{center}
  \begin{tikzpicture}[scale=1]
    \node[posladung] at (0,0) {$+$};
  \end{tikzpicture}
  \end{center}
  \vspace{2.1cm}
  \end{answerbox}
  \begin{solution}
  \Bild{5cm}{Bilder/punktladung_positiv.png}
  \end{solution}

  \part[1] Ist das Feld homogen oder inhomogen? Begründen Sie kurz.
  \begin{answerbox}
  \vspace{0.7cm}
  \end{answerbox}
  \begin{solution}
  Inhomogen: Die Feldlinien laufen radial auseinander. Mit wachsendem
  Abstand von der Ladung werden sie weniger dicht, das Feld wird also
  schwächer.
  \end{solution}
\end{parts}
\end{taskbox}
\end{questions}

% =============================================================================
% Aufgabe 2: Einzelne negative Ladung
% =============================================================================
\needspace{9.5cm}
\begin{questions}
\begin{taskbox}[Aufgabe 2: Einzelne negative Ladung]
\question[4]
\begin{parts}
  \part[3] Zeichnen Sie die Feldlinien um die abgebildete negative
  Ladung ein.
  \begin{answerbox}
  \vspace{3.4cm}
  \begin{center}
  \begin{tikzpicture}[scale=1]
    \node[negladung] at (0,0) {$-$};
  \end{tikzpicture}
  \end{center}
  \vspace{3.4cm}
  \end{answerbox}
  \begin{solution}
  \Bild{5cm}{Bilder/punktladung_negativ.png}
  \end{solution}

  \part[1] Ist das Feld homogen oder inhomogen? Begründen Sie kurz.
  \begin{answerbox}
  \vspace{0.7cm}
  \end{answerbox}
  \begin{solution}
  Inhomogen: Die Feldlinien laufen radial auf die Ladung zu. Mit
  wachsendem Abstand von der Ladung werden sie weniger dicht, das Feld
  wird also schwächer.
  \end{solution}
\end{parts}
\end{taskbox}
\end{questions}

% =============================================================================
% Aufgabe 3: Dipol
% =============================================================================
\needspace{13cm}
\begin{questions}
\begin{taskbox}[Aufgabe 3: Dipol]
\question[5]
\begin{parts}
  \part[4] Zeichnen Sie die Feldlinien zwischen den beiden abgebildeten
  Ladungen ein.
  \begin{answerbox}
  \vspace{3.7cm}
  \begin{center}
  \begin{tikzpicture}[scale=1.15]
    \node[posladung] at (-1.6,0) {$+$};
    \node[negladung] at (1.6,0) {$-$};
  \end{tikzpicture}
  \end{center}
  \vspace{3.7cm}
  \end{answerbox}
  \begin{solution}
  \Bild{7cm}{Bilder/dipol.png}
  \end{solution}

  \part[1] Ist das Feld homogen oder inhomogen? Begründen Sie kurz.
  \begin{answerbox}
  \vspace{0.7cm}
  \end{answerbox}
  \begin{solution}
  Inhomogen: Die Feldlinien sind gekrümmt, ihre Richtung und Dichte
  ändern sich von Ort zu Ort.
  \end{solution}
\end{parts}
\end{taskbox}
\end{questions}

% =============================================================================
% Aufgabe 4: Plattenkondensator
% =============================================================================
\needspace{11.5cm}
\begin{questions}
\begin{taskbox}[Aufgabe 4: Zwei parallele Platten]
\question[6]
\begin{parts}
  \part[4] Zeichnen Sie die Feldlinien zwischen und um die beiden
  abgebildeten, entgegengesetzt geladenen Platten ein.
  \begin{answerbox}
  \vspace{2.4cm}
  \begin{center}
  \begin{tikzpicture}[scale=1]
    \draw[very thick,red] (-3,1) -- (3,1);
    \foreach \x in {-2.6,-1.95,...,2.6} { \node[red] at (\x,1.35) {$+$}; }
    \draw[very thick,blue] (-3,-1) -- (3,-1);
    \foreach \x in {-2.6,-1.95,...,2.6} { \node[blue] at (\x,-1.35) {$-$}; }
  \end{tikzpicture}
  \end{center}
  \vspace{2.4cm}
  \end{answerbox}
  \begin{solution}
  \Bild{9cm}{Bilder/plattencondensator.png}
  \end{solution}

  \part[2] Ist das Feld homogen oder inhomogen? Begründen Sie kurz.
  \begin{answerbox}
  \vspace{1.1cm}
  \end{answerbox}
  \begin{solution}
  In der Mitte zwischen den Platten, fern vom Rand, homogen: Dort
  verlaufen die Feldlinien parallel und gleich dicht. Am Rand dagegen
  inhomogen, dort krümmen sich die Feldlinien nach aussen (Streufeld).
  \end{solution}
\end{parts}
\end{taskbox}
\end{questions}

% =============================================================================
% Abschluss: Van-de-Graaf-Generator
% =============================================================================
\begin{beispielbox}[Anwendung: der Van-de-Graaf-Generator]
Ein Förderband transportiert laufend Ladung zur Kugel oben hinauf, wodurch
sich diese immer stärker positiv auflädt. Eine zweite, geerdete Kugel
nähert sich: Durch das Feld der ersten Kugel werden aus der Erde
Elektronen angezogen, sodass sich die zweite Kugel negativ auflädt, ganz
ohne Berührung, genau wie einleitend besprochen. Wird das elektrische
Feld zwischen den beiden Kugeln stark genug, springt ein Funke über.

\Bild{7.5cm}{Bilder/van_de_Graaf_generator.png}
{\scriptsize Aufbau und Funktionsweise eines Van-de-Graaf-Generators.}
\end{beispielbox}

\end{document}
